\begin{lemma}[Trace decompositions from fixed-boundary restrictions]
    \label{lem:pgvwc07_trace_decompositions_from_block_restrictions}
    \lean{MPSTensor.pgvwc07_trace_decompositions_of_iSup_restrictions}
    \leanok
    Let $A^1,\ldots,A^g$ be block tensors. If
    $\psi\in(\C^d)^{\otimes(n+2)}$ satisfies
    \begin{align}
        \psi(a,-)&\in\bigvee_jG_{n+1}(A^j),
        \notag\\
        \psi(-,b)&\in\bigvee_jG_{n+1}(A^j)
        \notag
    \end{align}
    for every physical index $a,b$, then there are matrices $C^j_a$ and
    $D^j_b$ such that
    \begin{align}
        \psi&=\sum_j\alpha_j, \notag\\
        \alpha_j(i_1,\ldots,i_{n+2})
        &=
        \tr\!\left(A^j_{i_{n+2}}C^j_{i_1}
            A^j_{i_2}\cdots A^j_{i_{n+1}}\right), \notag
    \end{align}
    and, for every middle word $w$ of length $n$,
    \begin{align}
        \sum_j\tr\!\left(A^j_bC^j_aA^j_w\right)
        &=
        \sum_j\tr\!\left(D^j_bA^j_aA^j_w\right). \notag
    \end{align}
\end{lemma}

\begin{proof}
    \leanok
    Choose finite decompositions of the fixed-boundary vectors:
    \begin{align}
        \psi(a,i_2,\ldots,i_{n+2})
        &=
        \sum_j
        \tr\!\left(A^j_{i_2}\cdots A^j_{i_{n+2}}C^j_a\right), \notag\\
        \psi(i_1,\ldots,i_{n+1},b)
        &=
        \sum_j
        \tr\!\left(A^j_{i_1}\cdots A^j_{i_{n+1}}D^j_b\right). \notag
    \end{align}
    For every middle word $w$,
    \begin{align}
        \psi(a,w,b)
        &=
        \sum_j\tr\!\left(A^j_wA^j_bC^j_a\right) =
        \sum_j\tr\!\left(A^j_bC^j_aA^j_w\right), \notag
    \end{align}
    and the second decomposition gives
    \begin{align}
        \psi(a,w,b)
        &=
        \sum_j\tr\!\left(A^j_aA^j_wD^j_b\right) =
        \sum_j\tr\!\left(D^j_bA^j_aA^j_w\right). \notag
    \end{align}
    Hence
    \begin{align}
        \sum_j\tr\!\left(A^j_bC^j_aA^j_w\right)
        &=
        \sum_j\tr\!\left(D^j_bA^j_aA^j_w\right). \notag
    \end{align}
    For an arbitrary word $(i_1,\ldots,i_{n+2})$, the first decomposition
    gives
    \begin{align}
        \psi(i_1,\ldots,i_{n+2})
        &=
        \sum_j
        \tr\!\left(A^j_{i_2}\cdots A^j_{i_{n+2}}C^j_{i_1}\right) =
        \sum_j
        \tr\!\left(A^j_{i_{n+2}}C^j_{i_1}
            A^j_{i_2}\cdots A^j_{i_{n+1}}\right), \notag
    \end{align}
    by cyclicity of the trace.
\end{proof}

\begin{lemma}[One-step block intersection from block restrictions]
    \label{lem:block_restrictions_mem_block_ground_spaces}
    \lean{MPSTensor.pgvwc07_mem_iSup_groundSpace_of_iSup_restrictions}
    \leanok
    \uses{lem:pgvwc07_trace_decompositions_from_block_restrictions,
        lem:left_boundary_trace_decomposition_mem_block_ground_spaces}
    Let $A^1,\ldots,A^g$ be block tensors with common word-span separation at
    length $n$, and suppose that
    $\sum_a A^j_aA^{j\,\dagger}_a=\Id$ for every block $j$. If
    $\psi\in(\C^d)^{\otimes(n+2)}$ satisfies
    \begin{align}
        \psi(a,-)&\in\bigvee_jG_{n+1}(A^j),
        \notag\\
        \psi(-,b)&\in\bigvee_jG_{n+1}(A^j)
        \notag
    \end{align}
    for every physical index $a,b$, then
    $\psi\in\bigvee_jG_{n+2}(A^j)$.
\end{lemma}

\begin{proof}
    \leanok
    Choose finite decompositions of the fixed-boundary vectors:
    \begin{align}
        \psi(a,i_2,\ldots,i_{n+2})
        &=
        \sum_j
        \tr\!\left(A^j_{i_2}\cdots A^j_{i_{n+2}}C^j_a\right), \notag\\
        \psi(i_1,\ldots,i_{n+1},b)
        &=
        \sum_j
        \tr\!\left(A^j_{i_1}\cdots A^j_{i_{n+1}}D^j_b\right). \notag
    \end{align}
    Evaluating these two decompositions on $(a,w,b)$, where $w$ has length
    $n$, gives
    \begin{align}
        \sum_j\tr\!\left(A^j_bC^j_aA^j_w\right)
        &=
        \sum_j\tr\!\left(D^j_bA^j_aA^j_w\right). \notag
    \end{align}
    The first decomposition also gives
    \begin{align}
        \psi&=\sum_j\alpha_j, \notag\\
        \alpha_j(i_1,\ldots,i_{n+2})
        &=
        \tr\!\left(A^j_{i_{n+2}}C^j_{i_1}
            A^j_{i_2}\cdots A^j_{i_{n+1}}\right). \notag
    \end{align}
    Lemma~\ref{lem:left_boundary_trace_decomposition_mem_block_ground_spaces}
    therefore gives $\psi\in\bigvee_jG_{n+2}(A^j)$.
\end{proof}

\begin{lemma}[One-step block intersection characterization]
    \label{lem:block_restrictions_iff_block_ground_spaces}
    \lean{MPSTensor.pgvwc07_mem_iSup_groundSpace_iff_iSup_restrictions}
    \leanok
    \uses{lem:block_restrictions_mem_block_ground_spaces}
    Under the hypotheses of
    Lemma~\ref{lem:block_restrictions_mem_block_ground_spaces}, an
    $(n+2)$-site vector satisfies
    $\psi\in\bigvee_jG_{n+2}(A^j)$ if and only if
    \begin{align}
        \psi(a,-)&\in\bigvee_jG_{n+1}(A^j),
        \notag\\
        \psi(-,b)&\in\bigvee_jG_{n+1}(A^j)
        \notag
    \end{align}
    for every physical index $a,b$.
\end{lemma}

\begin{proof}
    \leanok
    The reverse implication is
    Lemma~\ref{lem:block_restrictions_mem_block_ground_spaces}. For the
    forward implication, write
    \begin{align}
        \psi&=\sum_j\gamma_j,
        \notag\\
        \gamma_j&\in G_{n+2}(A^j).
        \notag
    \end{align}
    Each $\gamma_j$ has fixed-boundary restrictions in $G_{n+1}(A^j)$:
    \begin{align}
        \gamma_j(a,-)&\in G_{n+1}(A^j),
        \notag\\
        \gamma_j(-,b)&\in G_{n+1}(A^j).
        \notag
    \end{align}
    Taking the finite sum over $j$ gives the two restrictions of $\psi$ in
    $\bigvee_jG_{n+1}(A^j)$.
\end{proof}

\begin{lemma}[Subspace form of the one-step block intersection]
    \label{lem:block_restriction_intersection_subspace}
    \lean{MPSTensor.pgvwc07_iSup_groundSpace_eq_restriction_intersection}
    \leanok
    \uses{lem:block_restrictions_iff_block_ground_spaces}
    Under the hypotheses of
    Lemma~\ref{lem:block_restrictions_mem_block_ground_spaces}, set
    $S_n:=\bigvee_jG_{n+1}(A^j)$. Then
    \begin{align}
        \left\{\psi:\forall b,\ \psi(-,b)\in S_n\right\}
        \cap
        \left\{\psi:\forall a,\ \psi(a,-)\in S_n\right\}
        &=
        \bigvee_jG_{n+2}(A^j). \notag
    \end{align}
\end{lemma}

\begin{proof}
    \leanok
    Lemma~\ref{lem:block_restrictions_iff_block_ground_spaces} gives the
    displayed equivalence for each vector $\psi$. Interpreting the two sides
    of that equivalence as subspaces gives the stated equality.
\end{proof}

\subsection{Eventual block-intersection consequences}
\label{subsec:block_intersection_eventual}

Propagating the single-length product span through a period window gives the
block-intersection identity for all sufficiently large $n$, both from a general
period window and from the normalized BNT representative hypotheses. The full
product span at length $n$ also makes the sum of the local spaces $G_n(A^j)$
direct.

\begin{lemma}[Period-window form of the block intersection identity]
    \label{lem:period_window_block_restriction_intersection}
    \lean{MPSTensor.pgvwc07_iSup_restriction_intersection_eventually_of_period_window}
    \leanok
    \uses{lem:period_window_word_tuple_span_padding,
        lem:block_restriction_intersection_subspace}
    Let $p>0$. Suppose the length-$p$ simultaneous block-word tuples span
    $\prod_j \MN{D_j}$, and suppose that, for every $0\le r<p$, the
    length-$(s+r)$ simultaneous block-word tuples span
    $\prod_j \MN{D_j}$. Assume also the normalization
    $\sum_a A^j_aA^{j\,\dagger}_a=\Id$ for every block $j$. Then there is a
    length $N$ such that, for every $n\ge N$, if
    $S_n:=\bigvee_jG_{n+1}(A^j)$, then
    \begin{align}
        \left\{\psi:\forall b,\ \psi(-,b)\in S_n\right\}
        \cap
        \left\{\psi:\forall a,\ \psi(a,-)\in S_n\right\}
        &=
        \bigvee_jG_{n+2}(A^j). \notag
    \end{align}
\end{lemma}

\begin{proof}
    \leanok
    Lemma~\ref{lem:period_window_word_tuple_span_padding} gives a length $N$
    such that the simultaneous block-word tuples span
    $\prod_j \MN{D_j}$ at every length $n\ge N$. For each such $n$,
    Lemma~\ref{lem:block_restriction_intersection_subspace} applies and gives
    the displayed subspace equality.
\end{proof}

\begin{lemma}[BNT block separation gives one block-intersection step]
    \label{lem:bnt_direct_sum_block_separation_block_intersection}
    \lean{MPSTensor.pgvwc07_iSup_restriction_intersection_of_bnt_directSum_selectors_c1}
    \leanok
    \uses{lem:bnt_direct_sum_block_separation_word_span,
        lem:block_restriction_intersection_subspace}
    Under the hypotheses of
    Lemma~\ref{lem:bnt_direct_sum_block_separation_word_span}, assume in
    addition the normalization
    $\sum_a A^j_aA^{j\,\dagger}_a=\Id$ for every block $j$. Let
    \begin{align}
        n&=(L_0+1)+(r-1)3(L_0+1),
        \notag\\
        S_n&=\bigvee_jG_{n+1}(A^j).
        \notag
    \end{align}
    Then
    \begin{align}
        \left\{\psi:\forall b,\ \psi(-,b)\in S_n\right\}
        \cap
        \left\{\psi:\forall a,\ \psi(a,-)\in S_n\right\}
        &=
        \bigvee_jG_{n+2}(A^j). \notag
    \end{align}
\end{lemma}

\begin{proof}
    \leanok
    Lemma~\ref{lem:bnt_direct_sum_block_separation_word_span} gives
    $\spn\{(A^1_w,\ldots,A^r_w): |w|=n\} = \prod_j\MN{D_j}$.
    Substituting this span into
    Lemma~\ref{lem:block_restriction_intersection_subspace}, together with the
    normalization, gives the stated subspace equality.
\end{proof}

\begin{lemma}[BNT block separation with a block-injectivity window]
    \label{lem:bnt_direct_sum_block_separation_period_window_block_intersection}
    \lean{MPSTensor.pgvwc07_iSup_restriction_intersection_eventually_of_bnt_directSum_period_window}
    \leanok
    \uses{lem:direct_sum_two_block_dimension_step,
        lem:pair_trace_separation_to_pairwise_block_separation,
        lem:block_separating_equations_from_pairwise_separation,
        lem:period_window_block_restriction_intersection}
    Let $A^1,\ldots,A^r$ be normalized representatives of a basis of normal
    tensors: each representative is irreducible, the family is left-canonical,
    the self-overlaps are normalized, and no two distinct same-dimensional
    representatives are gauge-phase equivalent. Assume $L>1$, and assume that,
    for every block $j$,
    \begin{align}
        \spn\{A^j_a:a=0,\ldots,d-1\}
        &=\MN{D_j}, \notag\\
        \spn\{A^j_u:|u|=L\}
        &=\MN{D_j}, \notag\\
        \spn\{A^j_v:|v|=3L\}
        &=\MN{D_j}. \notag
    \end{align}
    Set $S=3L$ and $q=(r-1)S$. Suppose that $p>0$ and let $s_0$ be a
    starting length. Assume that, for every block $j$,
    $\spn\{A^j_u:|u|=p\} =\MN{D_j}$,
    and that, for every $0\le s<p+q$ and every block $j$,
    $\spn\{A^j_v:|v|=s_0+s\} =\MN{D_j}$.
    Assume also the normalization
    $\sum_a A^j_aA^{j\,\dagger}_a=\Id$ for every block $j$. Then there is
    $N$ such that, for every $n\ge N$, if
    $S_n:=\bigvee_jG_{n+1}(A^j)$, then
    \begin{align}
        \left\{\psi:\forall b,\ \psi(-,b)\in S_n\right\}
        \cap
        \left\{\psi:\forall a,\ \psi(a,-)\in S_n\right\}
        &=
        \bigvee_jG_{n+2}(A^j). \notag
    \end{align}
\end{lemma}

\begin{proof}
    \leanok
    The direct-sum dimension step gives, for each ordered pair of distinct
    blocks $k\ne j$, coefficients $c^{k,j}_\omega$ of length $S$ such that
    \begin{align}
        \sum_{|\omega|=S} c^{k,j}_\omega A^k_\omega
        &=\Id, \notag\\
        \sum_{|\omega|=S} c^{k,j}_\omega A^j_\omega
        &=0. \notag
    \end{align}
    Multiplying these equations over the $r-1$ blocks $j\ne k$ gives
    coefficients $c^{(k)}_\omega$ of total length $q$ satisfying
    \begin{align}
        \sum_{|\omega|=q}c^{(k)}_\omega A^j_\omega
        &=
        \begin{cases}
            \Id, & j=k, \\
            0, & j\ne k.
        \end{cases} \notag
    \end{align}
    Combining this block-separating equation with block injectivity at length
    $p$ gives
    $\spn\{(A^1_w,\ldots,A^r_w):|w|=p+q\} = \prod_j \MN{D_j}$.
    Combining the same separating equations with the assumed
    block-injectivity window gives the full product span at lengths $s_0+q+s$,
    for $0\le s<p+q$. The period-window form of
    Lemma~\ref{lem:period_window_block_restriction_intersection} then gives
    the displayed equality for all sufficiently large $n$.
\end{proof}

\begin{lemma}[Unitality propagates simultaneous product spans]
    \label{lem:unital_simultaneous_product_span_propagation}
    \lean{MPSTensor.wordTupleSpanTop_succ_of_unital}
    \lean{MPSTensor.wordTupleSpanTop_of_ge_of_unital}
    \leanok
    Suppose that the length-$q$ simultaneous products span the product algebra,
    $\spn\{(A^1_w,\ldots,A^r_w):|w|=q\} = \prod_j\MN{D_j}$,
    and that every block satisfies
    $\sum_a A^j_aA^{j\,\dagger}_a=\Id$. Then the simultaneous products span
    the product algebra at every length $n\ge q$. This is the
    simultaneous-tuple form of the unital propagation in
    \cite[lines~893--898]{PerezGarcia2007Matrix}.
\end{lemma}

\begin{proof}
    \leanok
    \uses{lem:block_separating_equations_from_pairwise_separation}
    Given matrices $M_j$, first represent the tuple
    $((A^j_a)^\dagger M_j)_j$ by words of length $q$. Prefixing the letter $a$
    and summing over $a$ gives
    $\left(\sum_aA^j_aA^{j\,\dagger}_aM_j\right)_j =(M_j)_j$.
    This proves the step from $q$ to $q+1$; induction gives every $n\ge q$.
\end{proof}

\begin{lemma}[Normalized block-separating equations give all large product spans]
    \label{lem:unital_block_separating_product_span}
    \lean{MPSTensor.wordTupleSpanTop_of_ge_of_common_blockInjective_of_unital_of_pairBlockSeparatingWords}
    \lean{MPSTensor.wordTupleSpanTop_of_ge_of_bnt_directSum_unital_c1}
    \leanok
    \uses{def:pairwise_block_separating_words,
        lem:bnt_direct_sum_block_separation_word_span,
        thm:wordspan_propagation_unital}
    Let $A^1,\ldots,A^r$ be blocks satisfying the normalization
    $\sum_a A^j_aA^{j\,\dagger}_a=\Id$ and assume that every block is
    injective at length $L$:
    $\spn\{A^j_u:|u|=L\} =\MN{D_j}$.
    If there are block-separating equations of length $S$, then, for every
    $n\ge L+(r-1)S$,
    $\spn\{(A^1_w,\ldots,A^r_w):|w|=n\} = \prod_j \MN{D_j}$.
    In particular, for normalized BNT representatives satisfying the displayed
    unital normalization, if $L_0>0$ and $\Gamma_{L_0}^{A^j}$ is injective for
    every block $j$, the same conclusion holds with
    $S=3(L_0+1)$ and $n\ge (L_0+1)+(r-1)3(L_0+1)$.
\end{lemma}

\begin{proof}
    \leanok
    By Theorem~\ref{thm:wordspan_propagation_unital}, the normalization
    propagates the length-$L$ span equality to every larger homogeneous length:
    \begin{align}
        m\ge L
        \quad\Longrightarrow\quad
        \spn\{A^j_u:|u|=m\}
        &=\MN{D_j}. \notag
    \end{align}
    For $n\ge L+(r-1)S$, take $m=n-(r-1)S$. Combining the length-$m$ block
    span with the $r-1$ block-separating equations gives the displayed product
    span at length $n$. In the BNT case, the injectivity of
    $\Gamma_{L_0}^{A^j}$ gives the full span at length $L_0$, and the
    normalization propagates it to the lengths $L_0+1$ and $3(L_0+1)$.
    Lemma~\ref{lem:bnt_direct_sum_block_separation_word_span} then supplies the
    block-separating equations with $S=3(L_0+1)$.
\end{proof}

\begin{theorem}[Prepared-gauge simultaneous product span]
    \label{thm:pgvwc_prepared_gauge_product_span}
    \lean{MPSTensor.wordTupleSpanTop_threeBlock_mul_pred_of_blocksNotGaugePhaseEquiv_c1_of_dualFixedPoint}
    \leanok
    \uses{def:blockwise_product_word_span,
        def:blocks_not_gauge_phase_equiv}
    Let $r\ge2$, and let the bond dimensions $D_1,\ldots,D_r$ be positive.
    Suppose that every block $A^j$ is irreducible and that no two distinct
    blocks of the same bond dimension are gauge-phase equivalent. Let $L_0>0$,
    and suppose
    that, for every $j$,
    \begin{align}
        \spn\{A^j_w:|w|=L_0\}
        &=\MN{D_j}, \notag\\
        \spn\{A^j_w:|w|=L_0+1\}
        &=\MN{D_j}, \notag\\
        \spn\{A^j_w:|w|=3(L_0+1)\}
        &=\MN{D_j}. \notag
    \end{align}
    For every $j$, let $\Lambda^j$ be positive definite and satisfy
    \begin{align}
        \sum_a(A^j_a)^\dagger\Lambda^jA^j_a&=\Lambda^j. \notag
    \end{align}
    Then
    \begin{align}
        \spn\{(A^1_w,\ldots,A^r_w):|w|=3(r-1)(L_0+1)\}
        &=\prod_j\MN{D_j}. \notag
    \end{align}
    Diagonality of $\Lambda^j$ is not needed.
\end{theorem}

\begin{proof}
    \leanok
    \uses{lem:tp_gauge_equiv,
        thm:tp_gauge_from_adjoint_fixed,
        lem:gauge_invariance_block_injective,
        thm:is_normal_tensor_of_is_normal_left_canonical,
        thm:cpsv_normal_mpv_phase_gauge,
        lem:bnt_direct_sum_block_separation_word_span,
        thm:word_tuple_span_gauge_invariance}
    The square-root gauge is invertible by positive definiteness, and the dual
    fixed-point equation makes the prepared blocks left-canonical. Similarity
    preserves irreducibility, fixed-length injectivity, and pairwise inequivalence
    up to gauge and phase. Positive-length injectivity makes each prepared block
    algebraically normal; left-canonicality then makes it a normal tensor, whose
    self-overlap converges to one. Apply the existing sharp separation theorem to
    the prepared family and then use
    Theorem~\ref{thm:word_tuple_span_gauge_invariance} in the reverse direction.
\end{proof}

\begin{lemma}[Simultaneous product span at the source length]
    \label{lem:pgvwc_product_span_source_length}
    \lean{MPSTensor.wordTupleSpanTop_of_ge_of_bnt_directSum_unital_c1_pgvwc07_of_dualFixedPoint}
    \lean{MPSTensor.groundSpace_iSupIndep_of_ge_of_bnt_directSum_unital_c1_pgvwc07_of_dualFixedPoint}
    \leanok
    \uses{def:blockwise_product_word_span,
        def:blocks_not_gauge_phase_equiv,
        def:ground_space_parent}
    Let $r\ge2$, and let the bond dimensions $D_1,\ldots,D_r$ be positive.
    Suppose that every block $A^j$ is irreducible and that no two distinct
    blocks of the same bond dimension are gauge-phase equivalent. Let $L_0>0$,
    and suppose
    that every block is injective at length $L_0$:
    \begin{align}
        \spn\{A^j_w:|w|=L_0\}&=\MN{D_j}. \notag
    \end{align}
    For every $j$, let $\Lambda^j$ be positive definite and assume the PGVWC07
    canonical identities
    \begin{align}
        \sum_aA^j_aA^{j\,\dagger}_a&=\Id,
        &\sum_a(A^j_a)^\dagger\Lambda^jA^j_a&=\Lambda^j. \notag
    \end{align}
    For every $n\ge 3(r-1)(L_0+1)$, one has
    $\spn\{(A^1_w,\ldots,A^r_w):|w|=n\} =\prod_j\MN{D_j}$.
    Consequently the spaces $G_n(A^j)$ form an internal sum. No identity
    $\sum_a(A^j_a)^\dagger A^j_a=\Id$ is assumed for the source blocks.
    The doubly-normalized specialization remains available separately. This is
    the simultaneous spanning estimate used in
    \cite[Theorem~12, proof lines~1424--1456]{PerezGarcia2007Matrix}.
\end{lemma}

\begin{proof}
    \leanok
    \uses{lem:unital_simultaneous_product_span_propagation,
        thm:pgvwc_prepared_gauge_product_span,
        lem:product_span_block_image_direct_sum}
    Theorem~\ref{thm:pgvwc_prepared_gauge_product_span} gives the full
    simultaneous span at $q=3(r-1)(L_0+1)$ on the source representatives.
    The source unital identity and
    Lemma~\ref{lem:unital_simultaneous_product_span_propagation} propagate
    this equality to every $n\ge q$. To prove independence, write
    $\phi_j=\Gamma_n^{A^j}(X_j)$ and suppose that $\sum_j\phi_j=0$.
    Comparison of the coefficient of each word, followed by the simultaneous
    spanning equality, gives
    \begin{align}
        \sum_j \Gamma_n^{A^j}(X_j)=0
        &\quad\Longrightarrow\quad
        \left(\forall (M_j)_j,
            \sum_j\tr(X_jM_j)=0\right)
        \quad\Longrightarrow\quad
        \forall j,\ X_j=0. \notag
    \end{align}
    Indeed, in the last implication one takes a tuple supported in one block
    and uses nondegeneracy of the trace pairing. Hence every $\phi_j$ vanishes,
    so the spaces $G_n(A^j)$ form an internal sum.
\end{proof}

\begin{theorem}[Direct-sum restriction intersection at the source length]
    \label{thm:pgvwc_direct_sum_intersection_source_length}
    \lean{MPSTensor.pgvwc07_iSup_restriction_intersection_of_ge_of_bnt_directSum_unital_c1_pgvwc07_of_dualFixedPoint}
    \leanok
    \uses{def:ground_space_parent}
    Under the hypotheses of
    Lemma~\ref{lem:pgvwc_product_span_source_length}, let
    $m\ge3(r-1)(L_0+1)+1$. Then
    \begin{align}
        \C^d\otimes\left(\sum_jG_m(A^j)\right)
        \cap
        \left(\sum_jG_m(A^j)\right)\otimes\C^d
        &=
        \sum_jG_{m+1}(A^j). \notag
    \end{align}
    Lemma~\ref{lem:pgvwc_product_span_source_length} also shows that the sums
    at lengths $m$ and $m+1$ are internal. This is the direct-sum lemma in the
    proof of
    \cite[Theorem~12, proof lines~1346--1456]{PerezGarcia2007Matrix}.
\end{theorem}

\begin{proof}
    \leanok
    \uses{lem:pgvwc_product_span_source_length,
        lem:block_restriction_intersection_subspace}
    Let a vector belong to both restricted spaces. The two endpoint
    decompositions give boundary matrices $C_a^j$ and $D_b^j$. For every pair
    of physical indices $a,b$ and every middle word $w$ of length $m-1$, their
    coefficients obey
    \begin{align}
        \sum_j\tr\!\left((A_b^jC_a^j)A_w^j\right)
        &=
        \sum_j\tr\!\left((D_b^jA_a^j)A_w^j\right). \notag
    \end{align}
    The simultaneous product span at length $m-1$ and nondegeneracy of the
    trace pairing therefore give, block by block,
    $A_b^jC_a^j=D_b^jA_a^j$.
    Set $E^j=\sum_a C_a^jA_a^{j\,\dagger}$. The preceding identities and the
    right-unital normalization imply
    $A_b^jC_a^j=A_b^jE^jA_a^j$.
    Thus the component in block $j$ lies in $G_{m+1}(A^j)$. The reverse
    inclusion follows by splitting a word of length $m+1$ at either endpoint.
\end{proof}

\begin{remark}[Doubly normalized specialization]
    \label{rem:pgvwc_direct_sum_doubly_normalized_specialization}
    \lean{MPSTensor.wordTupleSpanTop_of_ge_of_bnt_directSum_unital_c1_pgvwc07}
    \lean{MPSTensor.groundSpace_iSupIndep_of_ge_of_bnt_directSum_unital_c1_pgvwc07}
    \lean{MPSTensor.pgvwc07_iSup_restriction_intersection_of_ge_of_bnt_directSum_unital_c1_pgvwc07}
    \leanok
    Under the hypotheses above, assume additionally that every block has
    normalized self-overlap and satisfies
    \begin{align}
        \sum_a(A^j_a)^\dagger A^j_a&=\Id. \notag
    \end{align}
    The retained doubly normalized declarations give the same simultaneous
    product span, internal-sum independence, and one-step
    restriction-intersection identity. They are the specialization
    $\Lambda^j=\Id$ of the source-normalized results.
\end{remark}

\begin{lemma}[Product span makes the sum of local spaces direct]
    \label{lem:product_span_block_image_direct_sum}
    \lean{MPSTensor.groundSpace_iSupIndep_of_wordTupleSpanTop}
    \lean{MPSTensor.groundSpace_iSupIndep_of_ge_of_bnt_directSum_unital_c1}
    \leanok
    \uses{def:ground_space_parent, lem:unital_block_separating_product_span}
    Suppose
    $\spn\{(A^1_w,\ldots,A^r_w):|w|=n\} = \prod_j\MN{D_j}$.
    Then the sum of the local spaces $G_n(A^j)$ is direct: if
    \begin{align}
        \phi_j&\in G_n(A^j),
        \notag\\
        \sum_j\phi_j&=0,
        \notag
    \end{align}
    then $\phi_j=0$ for every $j$. In particular, for the normalized BNT
    hypotheses above, this directness holds for every
    $n\ge (L_0+1)+(r-1)3(L_0+1)$.
\end{lemma}

\begin{proof}
    \leanok
    Write $\phi_j(\sigma)=\tr(A^j_\sigma X_j)$. Evaluating
    $\sum_j\phi_j=0$ on each word $w$ of length $n$ gives
    $\sum_j\tr(A^j_wX_j)=0$.
    Since $\tr(A^j_wX_j)=\tr(X_jA^j_w)$, the product-span equation gives, for
    every $(M_1,\ldots,M_r)\in\prod_j\MN{D_j}$,
    $\sum_j\tr(X_jM_j)=0$.
    Taking only the $j$-th component nonzero and using nondegeneracy of the
    trace pairing yields $X_j=0$, hence $\phi_j=0$. The BNT statement follows
    by substituting Lemma~\ref{lem:unital_block_separating_product_span}.
\end{proof}
